\documentclass[letterpaper,12pt,oneside]{article}
\usepackage[top=1in, left=1.25in, right=1.25in, bottom=1in]{geometry}

\usepackage[T1]{fontenc}
\usepackage[utf8]{inputenc}
\usepackage[spanish,es-nodecimaldot,es-tabla]{babel}
\usepackage{caption, subcaption}
\usepackage{graphicx}
\usepackage{array}
\usepackage{tikz}
\usepackage{imakeidx}
\usepackage[style=ieee]{biblatex}


\begin{filecontents*}{\jobname.bib}
@online{uapa_c,
  author  = {{Repositorio UAPA}},
  title   = {Contenido de C – Explicaciones y conceptos básicos},
  year    = {2026},
  url     = {https://repositorio-uapa.cuaed.unam.mx/repositorio/moodle/pluginfile.php/1023/mod_resource/content/2/contenido/index.html},
  urldate = {2026-02-16}
}

@online{ms_int,
  author  = {{Microsoft Docs}},
  title   = {Tipo int (Lenguaje C)},
  year    = {2026},
  url     = {https://learn.microsoft.com/es-es/cpp/c-language/type-int?view=msvc-170},
  urldate = {2026-02-16}
}

@online{c_perez,
  author  = {Carlos Pérez},
  title   = {Expresiones aritméticas en C},
  year    = {2026},
  url     = {https://www.carlospes.com/curso_de_lenguaje_c/01_08_01_expresiones_aritmeticas.php},
  urldate = {2026-02-16}
}

@online{ms_c_arit,
  author  = {{Microsoft Docs}},
  title   = {Operadores aritméticos (C\#)},
  year    = {2026},
  url     = {https://learn.microsoft.com/es-es/dotnet/csharp/language-reference/operators/arithmetic-operators},
  urldate = {2026-02-16}
}
\end{filecontents*}

\addbibresource{\jobname.bib}

\graphicspath{{./figs/}}
\usepackage{setspace}

\begin{document}

\begin{titlepage}
\setlength{\parindent}{0pt}
\setlength{\parskip}{0pt}

\begin{center}
\vfill

\includegraphics[width=4.5cm]{UNAM_LOGO2.png}

\vspace{0.5cm}
\Huge\textbf{Universidad Nacional Autónoma de México}

\vfill

\large Ingeniería en ...

\vfill
\large Estructura de Datos y Algoritmos I

\vspace{1cm}
{\huge FORMATO DE REPORTE}

\vspace{1cm}
ALUMNO:\\
\large 323197294

\bigskip
Grupo:\\
\#\\
Semestre:\\
2026-II

\vfill
México, CDMX. Febrero 2026

\end{center}
\end{titlepage}

\tableofcontents
\clearpage

\section{Introducción}

Este apartado debe abordar los siguientes puntos:

\begin{itemize}
\item \textbf{Planteamiento del problema.}
\item \textbf{Motivación.}
\item \textbf{Objetivos.}
\end{itemize}

\section{Marco Teórico}

En la programación estructurada, el lenguaje C es uno de los más importantes y utilizados debido a su eficiencia y versatilidad. 
El tipo de dato \texttt{int} es ampliamente utilizado para almacenar valores enteros dentro de los programas \cite{ms_int}. 
Las expresiones aritméticas permiten realizar operaciones matemáticas básicas como suma, resta, multiplicación y división \cite{c_perez}. 
Asimismo, los operadores aritméticos son esenciales para la manipulación de datos numéricos en distintos lenguajes de programación \cite{ms_c_arit}. 
Existen también recursos académicos que explican los fundamentos del lenguaje C y su aplicación práctica \cite{uapa_c}.

\section{Desarrollo}

Es la descripción de la implementación realizada en el lenguaje de programación, así como las pruebas realizadas para obtener los resultados.

\section{Resultados}

Mediante capturas de pantalla y una breve descripción seguida de la captura se presentan los resultados finales de su aplicación.

\section{Conclusiones}

Se presenta un análisis de los resultados obtenidos, donde se destaca la importancia de la aplicación de los conceptos teóricos para resolver el problema.

\printbibliography

\end{document}

